% !TeX program = pdflatex

% !TeX root = PF-11.tex

% Inventor (per express statement): Tero Keski-Valkama.

% Review route and source support in REVIEW_NOTES.md before filing.

\documentclass[pct]{sm-patent}

% Change pct to us for US Letter text. A4 drawings are separate.

\PatentTitle{Delay-mapped execution of selective temporal state with local dispatch}

\begin{document}

\descriptionstart

\pheading{Technical field}

\pnum Event-scheduled and physical-delay implementations of selective state-space processing.

\pheading{Background}

\pnum A selective recurrence can be evaluated by dense clocked arithmetic even when useful state access is sparse. Mapping its timing semantics to a local event implementation requires causal transport, address ownership and precision rules.

\pheading{Summary}

\pnum A processor maps input-dependent state transitions to scheduled elapsed-time transport, updates addressed state on delivered events, and dispatches outputs under declared timing and finite-precision rules.

\pheading{Brief description of the drawings}

\pnum FIG. 1 is a block diagram of an implementation. FIG. 2 shows operations of the implementation. Reference signs identify the same functional elements throughout the drawings.

\pheading{Detailed description}

\pnum In FIG. 1, 100 identifies events; 110 identifies state sram; 120 identifies selective maps; 130 identifies dispatcher; 140 identifies output; 150 identifies configuration. In FIG. 2, 210 identifies due time, 220 identifies arbitration and 230 identifies commit.

\pnum For a continuous-time diagonal mode with coefficient $A=-r+\mathrm{i}\omega$, an input-dependent step $\Delta_n\geq0$ gives $a_n=\exp(A\Delta_n)$. A write $b_n$ produces $z_n=a_nz_{n-1}+b_n$. An event embodiment stores $z$ and its last semantic timestamp and evaluates the elapsed transport only when the address is read or written.

\pnum A selective input map controls a delay, a write and a readout. Scheduler 130 assigns a modeled due time and dispatches only after causal prerequisites are available. A digital FPGA may represent due times as integer timestamps; it need not wait a wall-clock interval numerically equal to the modeled delay. A physical-delay circuit is a separate realization requiring its own calibrated delay transfer.

\pnum A fixed-point datapath stores real/imaginary pairs. Supplied or internally evaluated coefficients $c=e^{-r\Delta}\cos(\omega\Delta)$ and $s=e^{-r\Delta}\sin(\omega\Delta)$ implement $z'_R=cz_R-sz_I$ and $z'_I=sz_R+cz_I$. The design declares coefficient precision, rounding, saturation, timestamp wrap policy and protected nondecaying modes.

\pnum A race embodiment represents candidate clocks by completion times or clock values and selects an earliest eligible completion under a fixed tie rule. Keys and selected payload paths remain separate. Cancellation and replacement use versions or tokens so a stale timer cannot commit a new state write.

\pnum A precision-control embodiment applies common clock normalization under a fixed entering state and accounts for any transformation of race probabilities and delays. Route preservation does not alone establish equality of timing derivatives or behavior after state feedback. Clock noise and arithmetic quantization are characterized separately.

\pnum FIG. 1 shows input 100, local state 110, selective maps 120, event dispatcher 130, output 140 and configuration 150. FIG. 2 shows due-time computation 210, causal arbitration 220 and addressed commit 230. An ASIC implementation replaces FPGA memory and arithmetic mappings with technology-specific SRAM and standard cells while preserving the declared numerical/event contract. Area, timing and energy require implementation-specific measurement or labelled estimates.

\pheading{Hardware implementation}

\pnum A digital implementation may use synthesizable register-transfer logic mapped to an FPGA or to an application-specific integrated circuit. Addressed state uses register banks or SRAM macros; arithmetic uses integer, fixed-point or floating-point datapaths. Ready/valid or request/acknowledge interfaces apply backpressure without accepting a message twice. Synchronous FPGA execution, locally clocked circuits and asynchronous circuits are alternative implementations, each preserving its declared causal ordering and numerical precision.

\pnum Before deployment, a reference trace can check selected addresses, values, timestamps, publication versions and output probabilities against the circuit at the declared numerical precision. Timing quantization, ties, cancellation, overflow and saturation are specified for the implementation. FPGA resource utilization, post-layout timing and measured energy are implementation-dependent quantities; no particular measured hardware result is required by this embodiment.

\claimsstart

\begin{patentclaims}

\item A method of executing a selective temporal-state processor, comprising deriving an input-dependent elapsed-time parameter; evaluating transport of an addressed persistent state using a decay-and-rotation coefficient parameterized by the elapsed time; combining transported state with an input-dependent write; scheduling an output using a modeled completion time; and committing an addressed update only after its causal prerequisites are available.

\item The method of claim 1, wherein elapsed transport is evaluated upon access without periodic writes to unaddressed state.

\item The method of claim 1, wherein real and imaginary state coordinates are updated by fixed-point multiply-add operations with declared rounding and saturation.

\item The method of claim 1, wherein a digital scheduler represents modeled completion time independently of physical execution cycles.

\item The method of claim 1, wherein a version token rejects a cancelled or superseded timer before a state commit.

\item The method of claim 1, wherein candidate selection paths and selected payload paths are physically distinct.

\item An event-processing circuit comprising addressed memory, a temporal transport datapath, a due-time scheduler and a commit controller configured to perform the method of claim 1.

\item A non-transitory computer-readable storage medium storing instructions that cause a processor to perform the method of claim 1.

\end{patentclaims}

\abstractstart

A selective temporal processor stores addressed mode state and semantic timestamps. Input-dependent elapsed time controls lazy decay-and-rotation transport and a state write. A scheduler dispatches outputs after causal prerequisites are available and uses versions to reject stale timer actions. Digital timestamps can represent modeled delay independently of physical execution cycles. A fixed-point embodiment specifies coefficient precision, rounding and saturation. FPGA and ASIC implementations use the same declared event contract with technology-specific memory and arithmetic.

\end{document}

